\begin{abstract}
Bu çalışma, Tanzimat Fermanı ile birlikte yürürlüğe konulan ceza kanunnâmelerini merkeze alarak, fermanı takip eden yaklaşık yirmi yıllık süreçte kaleme alınan 1840 Kanunnâme-i Hümâyunu, 1851 Kanun-ı Cedîd ve 1858 Kanunnâme-i Hümâyunu arasında karşılaştırmalı bir inceleme yapmayı amaçlamaktadır. Fransız İhtilali sonrasında şekillenen vatandaşlık kavramı ve buna bağlı hak anlayışı, Tanzimat Fermanı'na da yansımış; ferman, tebaanın can, mal, ırz ve namus güvenliğini devletin teminatı altına alarak Osmanlı hukuk düzeninde vatandaşlık haklarının kabulüne yönelik önemli bir adım teşkil etmiştir.

Söz konusu güvenliğin hukuki zeminde tesis edilebilmesi amacıyla ilk olarak 1840 tarihli Kanunnâme-i Hümâyun kaleme alınmıştır. Ancak Tanzimat’la birlikte hız kazanan kurumsallaşma süreci, özellikle ilmiye sınıfı tarafından eleştirel ve zaman zaman muhalif bir tutumla karşılanmıştır. Bunun yanı sıra, uygulamada karşılaşılan sorunlar, 1840 tarihli ceza kanununun dönemin ihtiyaçlarına tam anlamıyla cevap veremediğini kısa sürede ortaya koymuştur. Bu durum, ceza hukukunun kapsamının genişletilmesi ve daha ayrıntılı bir düzenlemeye kavuşturulması gereğini doğurmuş; neticede 1851 tarihli Kanun-ı Cedîd hazırlanmıştır.

Öte yandan, Tanzimat sürecinin uluslararası boyutu da ceza hukuku alanındaki dönüşümü etkilemiştir. Kırım Harbi(1853–1856) ve ardından gerçekleştirilen Paris Barış Konferansı’nın(1856) yarattığı siyasi ve diplomatik atmosfer, 1858 tarihli Kanunnâme-i Hümâyun’un içeriğine de yansımıştır. Bu bağlamda çalışma, söz konusu üç ceza kanunnâmesinin muhteviyatını incelemekle birlikte, Tanzimat sonrası yirmi yıllık dönemde meydana gelen başlıca siyasi ve toplumsal gelişmeler ile ceza hukuku alanındaki düzenlemeler arasındaki ilişkiyi ortaya koymayı hedeflemektedir.

Çalışmada, bu süreçte Osmanlı ceza hukukunun giderek şer‘î hukuktan ayrılarak daha müstakil ve seküler bir yapıya doğru evrildiği ileri sürülmektedir. Bu dönüşümün temel nedenlerinden biri olarak, Tanzimat’la birlikte belirginleşen ilmiye–kalemiye çatışması gösterilmekte; söz konusu gerilimin ceza hukukunun normatif ve kurumsal çerçevesinin yeniden şekillenmesinde belirleyici bir rol oynadığı savunulmaktadır.

\textbf{Anahtar Kelimeler:} Şer`i Hukuk, Ceza Hukuku, Tanzimat Fermanı, Islahat Fermanı, Osmanlı Ceza Kanunnameleri
\end{abstract}

\begin{otherlanguage}{english} 
\begin{abstract}
This study focuses on the penal codes enacted in the wake of the Tanzimat Edict, offering a comparative analysis of the Ottoman Penal Codes of 1840 (Kanunname-i Hümayun), 1851 (Kanun-ı Cedid), and 1858 (Kanunname-i Hümayun), all of which were drafted within approximately twenty years following the promulgation of the edict. The concept of citizenship and the notion of citizens’ rights that emerged in the aftermath of the French Revolution were reflected in the Tanzimat Edict, which designated the state as the guarantor of the subjects’ security of life, property, honor, and sexual integrity, thereby laying the groundwork for the recognition of citizenship rights within the Ottoman legal order.

In order to establish this security on a legal basis, the first Ottoman penal code was promulgated in 1840. However, the institutionalization process introduced by the Tanzimat reforms was met with resistance, particularly from the members of the ilmiyye (religious-legal scholarly class). Moreover, it soon became apparent that the 1840 Penal Code was insufficient to meet the practical needs of the period. Consequently, the code was expanded and elaborated upon with the enactment of the 1851 Penal Code (Kanun-ı Cedid).

The transformation of Ottoman penal legislation during this period was also shaped by international developments. The Crimean War (1853–1856) and the subsequent Paris Peace Conference of 1856 exerted a significant influence on the formulation of the 1858 Ottoman Penal Code. Accordingly, this study not only examines the content of these three penal codes but also seeks to elucidate the relationship between the major political and social developments of the Tanzimat era and the evolution of Ottoman penal law over the course of these two decades.

The study argues that during this period Ottoman penal law gradually detached itself from sharīʿa-based legal frameworks and assumed a more autonomous legal character. It contends that this transformation was largely driven by the conflict between the ilmiyye and the kalemiyye (bureaucratic-administrative elite) that emerged alongside the Tanzimat reforms, a conflict that played a decisive role in reshaping the normative and institutional foundations of Ottoman penal law.

\textbf{Keywords:} Islamic Law, Penal Law, The Tanzimat Edict, The Islahat Edict, Ottoman Penal Legislation
\end{abstract}
\end{otherlanguage}
\clearpage